\paragraph{The coefficient tests at the prescribed cutoff.}
We keep the two level primes, the CM field \(E\) of degree \(2f\),
and the lift \(j<m\). Write \(\mathcal A\) for the sum of the
nonnegative flag exponents of the section being tested and \(n\)
for its central weight. For weight \(\kappa\) these numbers satisfy
\(\mathcal A+frn\le S_\kappa\). Let \(q_j\) be the negative signature rank. Define the
Fourier--Jacobi cutoff at this lift by
\[
 c_1=\tfrac12\operatorname{rd}(E)^{C_R}(2fq_j)^2q_j\,Tc,
 \qquad C_R=7 .
\]

The weight and cutoff fix the torsion order \(n'\) and the letter
radius bound \(\Lambda\) before the tests are chosen.
The split approximation uses weight allowance \(S_\kappa+r\),
level \(g_0g_1\), and the two integer norms
\(|\operatorname{disc}E|\) and \(\mathrm N\mathfrak c_K\).
The radius is enlarged to be at least one. The split construction
is used only at \(j<m\), where \(q_j\ge1\) and \(c_1>0\).
Thus the two signature cases entering the split construction are separated by
\[
  q_j=m-j\ge1\quad(j<m),\qquad q_m=0 .
\]
At \(j=m\) there is no complementary quotient.
\label{cor:slice-cm-coefficients-1}

An arc has a character lattice of rank one over \(\mathcal O_E\), a semiabelian model,
a normalized theta generator, and tangent letters in each
embedding block. Let \(t\) be its parameter and let \(X=(X_\tau)\)
be independent square frame matrices. Replace each letter in
row \(i\) by the combination with coefficients \((X_\tau)_{ik}\)
of the original letters.
Writing \(\ell_{\tau k}\) for the original letters, the variable frame has rows
\[
  \ell_{\tau i}(X)=\sum_k(X_\tau)_{ik}\,\ell_{\tau k}.
\]
 The section reading is a polynomial
\(P_F(X,t)\) in the entries of \(X\), with coefficients in
\(K'((t))\), over the number field of the arc. Homogeneity follows
by scaling all the tangent letters and applying the flag character.
In particular
\[
 P_F(X,t)\text{ has total frame degree }\mathcal A,
 \qquad \phi_{\mu,J}(F)=[X^\mu t^J]P_F(X,t).
\]
It is linear in \(F\). The same definition applies to a section
with complex scalar coefficients, by extending each coefficient
reading linearly along the chosen embedding of \(K'\).

\paragraph{Keeping the selected mode under a change of frame.}

To construct the CM arc, follow one nonzero Fourier coefficient through the change of flag.

Fix the component and a nonzero low coefficient \(F_\nu\) on a
rational isotropic line \(E\xi\). Let \(a_\tau(x)>0\) be the
cusp time at a flag frame \(x\), normalized as the reciprocal of
\(c_\tau\operatorname{maj}_{x,\tau}(\xi)\), with \(c_\tau\) the
positive flow constant of the component. The complex transvection
with parameter \(s=(s_\tau)_\tau\) sends these times to
\(a_\tau(x)+\operatorname{Im}s_\tau\). Its domain is the product
of half-planes given by positivity of the translated times.
Real transvections act on the Fourier coefficient by

\[
  \exp(4\pi i\sum_\tau\nu_\tau s_\tau).
\]
The product of the flowed coefficient and the reciprocal
exponential is holomorphic on these half-planes, by integrating
the section over its compact period torus. It is constant under
real translations. Its complex derivatives in those same
directions vanish, so the product is constant throughout the
domain. In particular transporting a nonzero reading to all
times equal to one keeps it nonzero.

For a flag write \(\widehat N(x)\) for the supremum of the
normalized point norms of \(F_\nu\) over its frames. This is
finite by the Fourier coefficient bound on the compact real
period torus, since those real transvections preserve the point
metric. It is positive at some flag of time one. Put

\[
  S=\sup_{a(x)=1}\widehat N(x)>0.
\]
 Take a flag of time one
whose supremum is greater than \(S/2\), then a frame on it
whose point norm is still greater than \(S/2\).
Signed Gram--Schmidt is a right upper triangular change of frame.
The weight character of the section and the norm factor cancel
under this change, and the majorant depends only on the signed
flag. We may therefore take an orthonormal frame \(x_0\) of the
same flag, with point norm between \(S/2\) and \(S\).

\begin{proposition}[A split frame detecting the selected mode]
\label{prop:split-frame-low-mode}
Fix the torsion order before the definite lift. If \(a_\nu\) is a nonzero selected low coefficient, it is detected by a bounded split approximation,
\[
 a_\nu\ne0\quad\Longrightarrow\quad \exists\,x_{\rm split}\quad a_\nu(x_{\rm split})\ne0.
\]

The construction keeps the component and isotropic line fixed.
It chooses the quotient complement and the frame after the coefficient
and order are fixed, using the metric and radius bounds below.
\end{proposition}

\begin{proof}
\proofstep{A motion that preserves a large coefficient.}
Put \(\rho=(12\pi(c_1+w_0))^{-1}\).
For a frame with times in
\([(1-\rho)/(1+\rho),(1+\rho)/(1-\rho)]\), move its times to one
as above. A bound for the ratio of the norm factors follows
from the leading Gram minors under that imaginary transvection.
The change in the Gram matrix under a complex transvection
with imaginary time \(s_\tau\) is a Hermitian rank-one
update. Its coefficient has the sign of \(s_\tau\).

For the \(i\)-th leading block let \(G_i(0)\) and \(G_i(s_\tau)\) be its Gram matrices before and after the update. The determinant lemma gives
\[
  \frac{\det G_i(s_\tau)}{\det G_i(0)}=1+s_\tau b_i.
\]
 Here \(b_i\) is the signed squared projection
of the line onto that block in the flow normalization.
Since the line is isotropic, the leading coefficients start
and end at zero, increase on the positive block and decrease
on the negative block, with maximum \(1/a_\tau\).
Thus the ratios lie between \(1\) and \(1+s_\tau/a_\tau\)
in either order.
More explicitly, with \(a_\tau=a_\tau(x)\), the norm factor at
\(x\) divided by the norm factor after transport is bounded by

\[
  \prod_\tau\max(a_\tau,a_\tau^{-1})^{s_\tau/2},
\]
 where
\(s_\tau\) is the spread of the weight at \(\tau\). Both time
ratios are at most \(1+3\rho\). This contributes at most
\(\exp(3\rho w_0/2)\le e\).
The quotient of the raw coefficient norms is
\(\exp(4\pi\sum_\tau \nu_\tau(1-a_\tau))\).
Since \(\nu_\tau\ge0\), \(\sum_\tau\nu_\tau\le c_1\) and
\(1-a_\tau\le2\rho\), the quotient is at most

\[
  \exp(8\pi\rho c_1)\le e.
\]
 Thus
\begin{equation}\label{eq:slice-split-near-sup}
  \widehat N(x)\le e^2 S .
\end{equation}
The movement and the comparison apply to every frame of the flag,
which explains the supremum on the left.

\paragraph{The first minimum on the isotropic quotient.}
Use \(L'\) for the Hermitian self-dual lattice of the component,
and \(\mathcal L\) for its period lattice. The fixed factors satisfy
\(N_A\mathcal L\subset L'\) and \(N_C L'\subset\mathcal L\). The
majorant-orthogonal complement of the hyperbolic plane at \(x_0\)
represents \(\xi^\perp/E\xi\). Write \(v_\tau\) for the image there
of a putative short nonzero class represented by \(v\in L'\cap\xi^\perp\).
Let \(N_0,H_0,K_0',R_0\) be the split parameters defined below, and
\(e_{\rm sp}=\min(1,f^2/(N_0^2 N_C^2 c_1 R_0))\).

We claim that
\[
 \max_\tau\operatorname{maj}_{x_0,\tau}(v_\tau)\ge e_{\rm sp}.
\]
If the maximum is smaller, choose \(0<\epsilon_0<e_{\rm sp}\)
bounding it. The Hermitian self-pairing of \(v\) is integral
and is the pairing of \(v_\tau\) at every embedding. It has
absolute value less than one everywhere, so it is zero.
Shift \(v\) by \(N_A b\xi\), \(b\in\mathcal O_E\), to make its
majorant projection coefficient \(b_\tau'\) onto \(\xi\) satisfy

\[
  |b_\tau'|\le 2f^2\operatorname{rd}(E)N_A
\]
 at all places.
Indeed the Euclidean dual of the embedded \(\mathcal O_E\)
is the conjugate codifferent scaled by two, by the trace formula.

For each nonzero element of the codifferent the field norm has
absolute value at least \(1/d_E\), by the inclusion of its
principal ideal in the codifferent. The arithmetic--geometric
mean inequality gives Euclidean dual minimum at least
\(2\sqrt f/\operatorname{rd}(E)\).
Let \(\rho_2\) be the Euclidean covering radius and \(\lambda_1^*\) the dual minimum. Banaszczyk's bound in dimension \(2f\) gives
\[
  \rho_2\le\frac{2f}{2\lambda_1^*}\le\frac{\sqrt f\,\operatorname{rd}(E)}2.
\]
The maximum norm radius is therefore at most
\(2f^2\operatorname{rd}(E)\). \cite{banaszczyk}
The shift leaves \(v_\tau\) and the assumption of shortness
unchanged. The two vectors \(\xi,N_Cv\) are now in \(\mathcal L\)
and span a totally isotropic plane.

\proofstep{The average along the transvection torus.}
Average along the real transvection torus for this plane. Its
characters are Hermitian two-by-two frequency matrices
\(\chi_\tau=\tau(\beta)\). The principal level torus contains the
translations obtained from integral Hermitian matrices after
multiplying by the level. Trace duality, including the off-diagonal
trace pairing, gives 
\[
  N_0\beta_{ij}\in\mathcal O_E,
\]
 with
\(N_0=(g_0g_1)d_E^2\). Taking the \(\xi\)-coefficient first retains
just \((\chi_\tau)_{00}=\nu_\tau\). The surviving frequencies are
positive definite at every embedding.

In fact translating in a
rank-one positive direction with negative frequency would give
exponential growth against the polynomially varying norm factor.
This gives positive semidefiniteness. If one determinant vanishes,
the algebraic determinant vanishes. A rational vector in the
kernel gives a rational isotropic line on which the frequency
is zero. The corresponding constant-term projection is zero by
cuspidality. Translation averages commute on this abelian torus,
so such a frequency has zero coefficient.

Normalize the direction \(v_\tau\) to have cusp time one, writing
it as \(a_{\tau,0}\xi+a_{\tau,1}N_Cv\) after this rescaling.
The normalization gives
\[
 |a_{\tau,1}|^2
   = (N_C^2 c_\tau\operatorname{maj}_{x_0,\tau}(v_\tau))^{-1},
 \qquad 0<c_\tau\le2 .
\]
Write \(d_\tau=\det\chi_\tau>0\), let \(R_\tau\) be the value of
this frequency on the normalized direction \(a_\tau\), and put
\(T=2\sum_\tau R_\tau\). Completing the square in the first
coordinate gives \(\nu_\tau R_\tau\ge d_\tau|a_{\tau,1}|^2\).

The algebraic integer \(\det(N_0\beta)\) is nonzero, so the
product of \(N_0^2d_\tau\) over the complex half is at least one.
Since \(\sum_\tau\nu_\tau\le c_1\), two applications of the
arithmetic--geometric mean inequality give
\[
 \sum_\tau d_\tau/\nu_\tau\ge \frac{f^2}{N_0^2c_1},
 \qquad T\ge\frac{f^2}{N_0^2N_C^2c_1\epsilon_0}> R_0 .
\]
The first inequality uses
\(\prod_\tau(f\nu_\tau/c_1)\le1\). In particular all the
\(\nu_\tau\) under consideration here are positive.

\proofstep{The first minimum at the same frequency.}
The height count is on the same frequencies. In each embedding
write the second standard vector as \(q a_\tau+p(1,0)\).
Then \[
 |q|^2=N_C^2 c_\tau\operatorname{maj}(v_\tau)\le2N_C^2,\qquad |p|^2=N_C^2|b_\tau'|^2.
\]
Positivity and Cauchy--Schwarz bound the second diagonal and
the off-diagonal frequency entries by \(H_0(c_1+T/2)\).
For the off-diagonal use

\[
  |(\chi_\tau)_{01}|^2\le\nu_\tau(\chi_\tau)_{11}.
\]
The first diagonal is prescribed, so the other two entries
determine the frequency. In the maximum of the complex norms
the integer lattice has minimum at least one, by the field norm.
Disjoint translates of the open ball of radius \(1/2\) give at
most \((1+2B)^{2f}\) lattice points in a ball of radius \(B\).
For \(t\ge1\) use \(B=N_0 H_0(c_1+t/2)\) for each of the two
entries. Counting both entries therefore bounds the number of Hermitian frequency matrices \(\beta\) by
\[
  \#\{\beta:T(\beta)\le t\}\le(K'_0t)^{4f}\qquad(t\ge1).
\]

The normalized \(v_\tau\)-transvection commutes with the
\(\xi\)-flow. The time-one flag, and every frame in it, can be
moved back by \(i/2\) in the normalized \(v_\tau\) direction
at each place. Its time becomes \(1/2\), and the \(\xi\)-time
stays one, since the two directions are majorant-orthogonal.
The real period averages of the rank-two group at that moved
point therefore bound the normalized coefficient by \(S\).
Transport back multiplies the raw coefficient by \(e^{-\pi T}\)
and changes the norm factor by at most \(2^{w_0/2}\).
Group the resulting bounds by \(m=\lceil T\rceil\). For
\(m>R_0=w_0+4f\log(8fK_0')+4f+2\), the inequality
\(\log(m/(8f))\le m/(8f)-1\) gives

\[
  (K_0'm)^{4f}\le e^{2m-w_0-2}.
\]
Thus the sum of the coefficient bounds, divided by \(S\), is
at most
\[
 2^{w_0/2} e^{\pi-w_0-2}
 \frac{e^{4(2-\pi)}}{1-e^{2-\pi}}<\tfrac12 .
\]
The absolute bound justifies recovering the section from its
Fourier series on the compact torus. Taking the supremum over
the original flag contradicts the choice above \(S/2\). This
proves the quotient minimum with the cutoff fixed at \(c_1\).
\label{cor:slice-split-frame-1}

\paragraph{A signed rational complement.}

The quotient on the isotropic line controls a first minimum. To turn that estimate into a CM frame we need a rational complement with the prescribed signs.

The quotient of \(L'\cap\xi^\perp\) by \(L'\cap E\xi\) is
a self-dual Hermitian lattice on \(\xi^\perp/E\xi\).
One can check this locally. A generator of the saturated line
is primitive, and self-duality supplies a partner pairing to one.
If these vectors are \(\xi_0,\eta_0\), the Gram
matrix on their plane is
\[
 \begin{pmatrix}0&1\\1&H(\eta_0,\eta_0)\end{pmatrix}
\]
and its determinant is \(-1\). Its inverse is integral,
so subtracting the plane projection preserves the lattice.
Inside this plane the perpendicular to \(\xi_0\) is
the line it spans. Thus the quotient lattice is the
integral orthogonal complement. Its signature in the selected convention
is \((q_j-1,r-q_j-1)\).

\proofstep{Choosing the rational complement by transference.}
Use the transference bound for this lattice with lower majorant
\(e_{\rm sp}\). In real dimension \(2f(r-2)\) it gives covering
radius and an independent short family at

\[
  \mu=4(f(r-2))^3d_E/\sqrt{e_{\rm sp}}.
\]
Use the same lattice construction as for the full interior frame
with approximation allowance \(\rho/40\). Namely, approximate
the positive axes scaled by \(4\sqrt{r-2}\mu/(\rho/40)\),
and multiply the orthogonal projections of the short spanning
family to the negative space by the positive Gram determinant.
The positive block errors have operator norm at most
\(\sqrt{r-2}\mu\), so its span is a graph of norm at most
\(\rho/120\). The negative projection is given by the Gram
adjugate and is integral.

The leading positive Gram determinant
and the Gram-adjugate formula give the length bound
\(B_{\rm rad}^{2q_j-1}\) on the resulting letters, with
\(B_{\rm rad}=2\cdot4\sqrt{r-2}\mu/(\rho/40)\).
Indeed, for a vector \((u,v)\) the negative projected norm
squared is
\[
 \|(1-AA^*)^{-1/2}(v-Au)\|^2\le (1+\|A\|^2)\|(u,v)\|^2/(1-\|A\|^2)\le2\|(u,v)\|^2.
\]
Raising the common bound
for the positive lengths to twice the number of positive
axes pays the determinant. If \(q_j=1\) there is no positive
axis in the quotient, and the short negative family needs
no projection. Lift the letters to \(L'\cap\xi^\perp\).
Their pairings and their quotient basis are independent
of the choices of lifts.

In the signed orthonormal basis put
\(C_+=(1-A^*A)^{-1/2}\) and \(C_-=(1-AA^*)^{-1/2}\).
The block matrix with the required positive graph is
\[
                 U_A=\begin{pmatrix}C_+&A^*C_-\\
                                        AC_+&C_-\end{pmatrix}.
\]
For the form with matrix \(\operatorname{diag}(1,-1)\),
the two column blocks are orthogonal: their mixed matrix
is \(C_+(A^*-A^*)C_-=0\). Their squared Gram blocks are
\[
 C_+(1-A^*A)C_+=1,\qquad C_-(AA^*-1)C_-=-1.
\]
Thus \(U_A\) is unitary. Extend it by the identity on the
hyperbolic plane.

The
resulting frame \(x_1\) has the selected letters in the two
eigenspaces on the complement, and its \(\xi\)-time is still
one. The positive plane in the chart of \(x_0\) has graph
\(Z_1\) of norm at most \(\rho/100\). With
\(t=(1-\rho/100)/(1+\rho/100)\), the majorant at \(x_1\) is
between \(t\) and \(1/t\) times that at \(x_0\). The quotient
lower bound for \(x_1\) is therefore \(t e_{\rm sp}\).

\paragraph{Clearing the projection at the prescribed order.}
Let \(W_\tau\) denote the complex complement at \(x_1\),
and write \(P_\tau(\mu)\) for the coefficient of \(\xi\)
in the projection of \(\mu\in\mathcal L\cap\xi^\perp\).
The coefficient projection comes from the exact sequence
\[
 0\longrightarrow \mathcal L\cap E\xi
 \longrightarrow \mathcal L\cap \xi^\perp
 \longrightarrow
   (\mathcal L\cap\xi^\perp)/(\mathcal L\cap E\xi)
 \longrightarrow 0.
\]
The quotient is torsion-free because the submodule is the
intersection with an \(E\)-line. It is finitely generated over
the Dedekind ring \(\mathcal O_E\), hence projective, so the
sequence splits over \(\mathcal O_E\). Write \(l(\mu)\xi\) for an integral
projection from such a splitting. Its coefficient extends
by scalar extension to an \(E\)-linear form with 
\[
  l(\xi)=1.
\]
Hence the difference \(P_\tau-l_\tau\) is a linear form on
the quotient.

\proofstep{Simultaneous clearing at the prescribed torsion order.}
We approximate \(n'(P_\tau-l_\tau)\) simultaneously by the
linear forms of one lattice element. In the real tangent
space \(\prod_\tau W_\tau\) with the Euclidean majorant metric,
the image of \(\mathcal L\cap\xi^\perp\) is a lattice of
dimension \(d=2f(r-2)\). The quotient minimum is at least
\(\sqrt{t e_{\rm sp}}/N_A\): multiplication by \(N_A\) sends
the projected lattice to the one used for the first-minimum bound.
The marked Riemann pairing is
\(\operatorname{Tr}_{E/\mathbb Q}(-\xi_0 H(-,-))\) on period
vectors, where 
\[
  2|\tau(\xi_0)|=k_0.
\]
This pairing is perfect on the projected lattice, since it is
perfect on the period lattice and the isotropic sublattice is
saturated.

To see this, extend a basis of the primitive
isotropic sublattice by symplectic partners, obtained by lifting
their dual functionals under the perfect pairing. The resulting
even-rank block is unimodular, and the perpendicular quotient
is its integral orthogonal complement.
Relative to the majorant, the real operator associated to the
pairing is \(k_0\) times an orthogonal map.
Writing \(T=k_0 U\) with \(U\) orthogonal for the majorant gives
\[
  \|Tx\|=k_0\|x\|,\qquad T^*T=k_0^2I .
\]
 Indeed it is the
product of multiplication by a purely imaginary scalar,
the signature involution, and the factor two from the CM trace.

The similarity \(T\) takes the projected lattice to its Euclidean dual.
Thus the dual minimum is at least \(k_0\sqrt{t e_{\rm sp}}/N_A\).
Transference and rounding in \(d\) independent vectors give
covering radius at most
\(d^3N_A/(2k_0\sqrt{t e_{\rm sp}})\) for the projected lattice.
Via Riesz representation, the linear forms
\[
 u\longmapsto d_E(-\tau(\xi_0)) H_\tau(u,\mu_{0,\tau}),
 \qquad \mu_0\in\mathcal L\cap\xi^\perp ,
\]
come from the projected lattice by an operator of norm
\(d_Ek_0/2\) times an isometry. The covering estimate
therefore approximates the tuple \(n'(P_\tau-l_\tau)\)
with error in its representing vector bounded by
\(d_Ed^3N_A/(4\sqrt{t e_{\rm sp}})\).
The majorant of \(\xi\) is \(1/k_0\).

Multiplying the error
norm by \(1/\sqrt{k_0}\) for its value on the line is
covered by \(B_{\rm ind}(t e_{\rm sp})\), the index allowance
below evaluated with lower bound \(t e_{\rm sp}\).
For the discriminant factor, the pairings \(H(L',L')\)
generate the unit ideal by Hermitian self-duality.
Since 
\[
  N_C L'\subset\mathcal L,
\]
 integrality of the trace
pairing puts \(N_C^2(-\xi_0)\) in the codifferent. All its
embedding norms equal \(N_C^2 k_0/2\). The ideal norm lower
bound \(1/d_E\) therefore gives
\(N_C^2 k_0\operatorname{rd}(E)\ge2\), as used in the index
allowance.
All these comparisons concern the projected lattice of the
full period module, independent of the particular spanning
letters just chosen.

\proofstep{Returning to the Hermitian frame.}
Dividing by \(n'\) gives a correction \(\delta_\tau\) on \(W_\tau\)
of operator norm, as a map into the line with its majorant metric,
at most \(\kappa=B_{\rm ind}(t e_{\rm sp})/n'\).
It has the required integral effect. For \(\mu\in\mathcal L\cap\xi^\perp\)
the corrected projection multiplied by \(n'\) has coefficient
\[
 n'l(\mu)+d_E(-\xi_0)H(\mu,\mu_0).
\]
The second summand is integral since its value before multiplication
by \(d_E\), paired against each \(a\in\mathcal O_E\) by the trace,
is \(\operatorname{Tr}_{E/\mathbb Q}(-\xi_0 H(a\mu,\mu_0))\in\mathbb Z\).
Thus it is in the trace dual. The first summand belongs to the
line coefficient module by the integral splitting. The vector
with the displayed coefficient is in \(\mathcal L\cap E\xi\).

At each embedding choose \(\xi'\) in the hyperbolic plane with
\(H_\tau(\xi,\xi')=1\), of norm squared the reciprocal of
\(\operatorname{maj}_{x_1,\tau}(\xi)\), and isotropic.
Take \(w_\tau\in W_\tau\) with \(H_\tau(u,w_\tau)=\delta_\tau(u)\).
The unitary fixing \(\xi\) acts by
\[
 u\longmapsto u+\delta_\tau(u)\xi,\qquad
 \xi'\longmapsto
 \xi'-w_\tau-\tfrac12 H_\tau(w_\tau,w_\tau)\xi .
\]
Preservation of all pairings follows by expanding them on the
three summands. In this moved frame the line projection is
\(P_\tau-\delta_\tau\) and the eigenspaces on the quotient are
unchanged. The two linear off-diagonal terms have operator
norm at most \(\kappa\) each and the quadratic term at most
\(\kappa^2/2\), so the move differs from the identity by at
most \(3\kappa\) in the \(x_1\) majorant when \(\kappa\le1\).

In the \(x_0\) majorant the bound is \(\eta=3\kappa/t\).
With \(z_1=\rho/100\), the transformed graph has norm at most
\[
 \frac{z_1+\eta(1+z_1)}{1-\eta(1+z_1)}<\rho/21 .
\]
Indeed the new upper block on \(\binom{1}{Z_1}\) differs from
the identity by at most \(\eta(1+z_1)\), and the lower block
changes by at most the same amount. The inequality on \(n'\)
gives 
\[
  \kappa\le\rho/(1000t),
\]
 and \(\rho\le1/10\) gives the
strict numerical bound. This construction uses the \emph{chosen}
positive integer \(n'\). Clearing denominators first and then
taking an arbitrary multiple would give a different quantifier.
\label{cor:slice-split-frame-2}

\paragraph{The small-chart nonvanishing test.}
Let \(Z\) be the lower-block chart parameter of a final moved
flag in the orthonormal frame \(x_0\). Use the operator norm on
each block, and the maximum over the complex places. For
\(\|W\|\le\rho\) the points \(x(W)\) of this chart are in the
domain. For each singular value \(s\) of \(W\), the majorant in the
associated signed two-plane has eigenvalues \((1+s)/(1-s)\)
and its reciprocal relative to \(x_0\).
The two relative eigenvalues are reciprocal, as can be read from
\[
  \frac{(1+s)^2}{1-s^2}=\frac{1+s}{1-s},\qquad \frac{(1-s)^2}{1-s^2}=\frac{1-s}{1+s}.
\]
 This follows from the
unitary graph matrix above, whose two-by-two block is
\((1-s^2)^{-1/2}\left(\begin{smallmatrix}1&s\\s&1\end{smallmatrix}\right)\).
The remaining eigenvalues are one. Thus the majorants change by
factors between \((1-\rho)/(1+\rho)\) and its reciprocal.

For a dominant weight
of spread at most \(w_0\), the norm factor \(J(x)\) satisfies
\[
  |\log J(x(W))-\log J(x_0)|\le 4w_0 r\|W\|^2\le \tfrac14 .
\]
For leading minors ending in the positive block one takes
determinants of \(1-W_1^*W_1\), with \(W_1\) the relevant columns
of \(W\). For larger minors, eliminate the retained negative
identity block. The remaining factor is \(1-W_2^*W_2\), with
\(W_2\) the rows of \(W\) after deleting that block. The full
determinant has absolute value one.

For each of the other positive absolute determinant ratios \(q_i\), the block bound gives
\[
  |\log q_i|\le-r\log(1-\|W\|^2).
\]
Summation by parts in the consecutive minor ratios charges
only the weight spread, giving the stated quadratic estimate.
This and the majorant comparison put the chart in the range of
\eqref{eq:slice-split-near-sup}. If \(\|Z\|<\rho/21\), put
\(R=\rho/\max(\|Z\|,\rho/22)>21\).
The function \(g(t)=J(x_0)F_\nu(x(tZ))\) on \(|t|\le R\)
has \(|g(t)|\le e^{9/4}S\), and \(|g(0)|>S/2\).
If \(g(1)=0\), precompose with the automorphism of this disc
that exchanges 0 and 1. Schwarz's lemma gives

\[
  |g(0)|\le e^{9/4}S/R<S/2,
\]
 a contradiction.
This is the form of the nonvanishing argument used in the split
detector. It allows the frame move to depend on the coefficient
and the prescribed torsion order.
\label{cor:slice-split-frame-3}
\end{proof}


\paragraph{Isolation of the frequency on one arc.}

Nonvanishing on a small chart still permits several Fourier frequencies. An arc on which the chosen frequency can be separated gives a single coefficient to estimate.

Let \(E^+\) be the fixed field of complex conjugation and
write \(\operatorname{Tr}=\operatorname{Tr}_{E/\mathbb Q}\).
Let \(N\) be the transvection level, \(Y\) the rank-one
character lattice and \(\xi:Y\to\mathbb Z\) its marked
nonzero functional. Choose \(D>0\) clearing the trace dual
of \(\mathcal O_E\).

\begin{lemma}[The frequency lattice at a split cusp]\label{lem:slice-cm-frequency-lattice}
Every occurring frequency \(\nu\) at the selected split cusp
satisfies the archimedean positivity and finite trace conditions
\[
 \sigma(\nu)\ge0\quad(\sigma:E^+\hookrightarrow\mathbb R),\qquad \operatorname{Tr}_{E^+/\mathbb Q}(ND\nu\,\mathcal O_{E^+})\subseteq\mathbb Z.
\]

\end{lemma}

\begin{proof}
The principal-level transvections for the selected isotropic
line have trivial weight character. The full transvection
period lattice is rational in \(E^+\). Thus Fourier characters
are trace pairings with elements of \(E^+\). For a nonzero
coefficient, sending one cusp time to infinity with the other
times fixed multiplies its raw term by the corresponding Fourier
exponential. For a nonzero coefficient \(a_\nu\), with the
other times fixed, the modulus of that Fourier reading is a
constant multiple of
\(|a_\nu|\exp(-2\pi t\,\tau(\nu))\).
It is at most the supremum on the real period torus.
The norm change on this time family grows at most polynomially
in \(t\) in both directions. If \(\tau(\nu)<0\) the displayed
exponential violates that bound as \(t\to\infty\).
This proves total nonnegativity.

Trace duality gives
\(p:Y\to E\) with 
\[
  \operatorname{Tr}(a p(y))=\xi(ay)
\]
 for
\(a\in\mathcal O_E\). The map is \(\mathcal O_E\)-linear and
injective since it is nonzero on the torsion-free rank-one
lattice. For \(t_0\in\mathcal O_{E^+}\) the
bilinear form
\[
  (y,y')\longmapsto
  \operatorname{Tr}\bigl(2NDt_0 p(y)\overline{p(y')}\bigr)
\]
is divisible by \(N\), since \(D\overline{p(y')}\) is integral.
Under the unimodular alternating period pairing this divisibility
says that the transvection with parameter \(NDt_0\) moves each
lattice vector by an element of \(N\) times the lattice. One sees
this by pairing its displacement against every lattice vector.
The transvection is therefore a period with trivial weight
character. Its Fourier character gives
\(\exp(2\pi i\operatorname{Tr}(ND\nu t_0))=1\) for every \(t_0\)
in the integer ring of \(E^+\), which is the asserted dual
lattice condition.
\end{proof}

\begin{proposition}[Isolating a low mode by arc coefficients]
\label{prop:cm-low-mode-detector}
For \(j<m\), let \(a_\nu\) denote extraction of the Fourier--Jacobi coefficient of frequency \(\nu\). On the selected source space, the allowed arc coefficient tests detect precisely the modes of half-trace at most \(c_1\),
\[
 \bigcap_{\mu,J,\ {\rm allowed\ arcs}}\ker\phi_{\mu,J}=\bigcap_{\nu:\,\sum_\tau\nu_\tau\le c_1}\ker a_\nu.
\]

Here \(\tau\) ranges over the chosen CM half, so
\(\sum_\tau\nu_\tau=\tfrac12\operatorname{Tr}_{E/\mathbb Q}(\nu)\).
The statement holds after extension along the chosen complex
embedding because each nonzero low coefficient has a detecting arc
and frame. The detecting field may depend on that coefficient.
\end{proposition}

\begin{proof}
\proofstep{Choosing the frequency on the arc.}
Put \(j_0=\operatorname{Tr}(ND\nu_0)\), a nonnegative integer.
The totally nonnegative elements of this trace dual of trace
\(j_0\) are finite. Their real conjugates lie in a compact
simplex and the dual embeds as a lattice. Choose
\(b\in\mathcal O_{E^+}\) such that
\(\operatorname{Tr}((\lambda-ND\nu_0)b)\ne0\) for all other
elements of the slice. The nondegenerate trace pairing makes
each excluded set a proper rational hyperplane, and an
integral lattice cannot be covered by finitely many such
hyperplanes. Set \(M=\max_\tau |\tau(b)|\) and take an
integer \(R>M(2j_0+1)\). The element \(u=R+b\) is totally
positive.

If 
\[
  \operatorname{Tr}(\lambda u)=
\operatorname{Tr}(ND\nu_0 u)
\]
 and \(i=\operatorname{Tr}\lambda\),
then \(R|i-j_0|\le M(i+j_0)\). If \(i\ne j_0\) then \(|i-j_0|\ge1\) and
\(i+j_0\le |i-j_0|+2j_0\). Combining the two inequalities gives
\[
 R|i-j_0|\le M(i+j_0)\le M(2j_0+1)|i-j_0|,
\]
contrary to the choice of \(R\).
Equality of the perturbed traces in the slice now forces
\(\lambda=ND\nu_0\).

\proofstep{The period and its Fourier variable.}
Use \(s_E=NDu\) in the period form above. This form is
positive definite and level-divisible. Conversely it determines
\(s_E\) uniquely by trace duality against the multiples of
one nonzero \(p(y)\overline{p(y)}\). The rank-one toroidal
uniformization therefore gives an arc with the times
\(y\tau(s_E)\) at height \(y\). Its exponent for \(\nu_0\) is
\[
 J=2\sum_{\tau\text{ in the CM half}}
              \tau(\nu_0)\tau(s_E).
\]
All occurring frequencies have integral exponents by the level
condition, and \(\nu_0\) is the unique one with exponent \(J\).
Along the deep part of the flow the Fourier expansion converges
normally. Integrating the circle of \(t=e^{2\pi i w}\) against
\(t^{-J}\) therefore selects exactly the \(\nu_0\)-term.

For invertible \(X\), the resulting identity has the form
\begin{equation}\label{eq:slice-arc-coeff}
     [t^J]P_F(X,t)\,t^J
          = \kappa_X\,F_{\nu_0}(x_X),
          \qquad \kappa_X\ne0 .
\end{equation}
Here \(F_{\nu_0}\) is measured at the corresponding complex
frame and height, and \(\kappa_X\) includes the change of
theta section and tangent frame. The normalized base change
identifies the arc theta generator with the complex generator
up to a nonzero scalar. Formula \eqref{eq:slice-arc-coeff}
holds on the complex chart and hence computes its Laurent
coefficient, even if there is a pole at the boundary. Rescaling
\(t\) by a nonzero constant changes both sides of the coefficient
comparison by the matching power.
\label{cor:slice-cm-coefficients-2}

\proofstep{The kernel of the coefficient tests.}
These identities give both directions of the kernel assertion.
Every allowed arc has
\(\tfrac12\operatorname{Tr}_{E/\mathbb Q}\nu_0\le c_1\).
If every mode of half-trace at most \(c_1\) vanishes, write \(Q_J(X)=[t^J]P_F(X,t)\). On the punctured arc, \eqref{eq:slice-arc-coeff} gives
\[
 Q_J(X)=t^{-J}\kappa_X F_{\nu_0}(x_X)=0
       \qquad (X\ \text{invertible}).
\]
The product of the general linear groups is Zariski dense in the product of the matrix spaces. Hence the polynomial \(Q_J\) vanishes identically, which kills every coefficient test on this arc. Conversely take a nonzero
coefficient on one component with

\[
  \sum_\tau\nu_\tau=\tfrac12\operatorname{Tr}\nu\le c_1.
\]
The split construction just proved keeps this coefficient nonzero
and gives the bounded CM basis on the complement of the hyperbolic
plane. Use the resulting frame and the above positive form
to construct the arc. The abelian period lattice is obtained
from the orthogonal to the isotropic line by quotienting by
the line. Its complex structure is diagonal on the split CM
letters. The principal Riemann form on the original lattice
identifies the rank-one quotient with the dual of the isotropic
lattice.
Writing \(I\) for the isotropic lattice and \(L\) for the ambient period lattice, this is the perfect duality induced by the principal alternating form:
\[
  L/I^\perp\simeq I^\vee .
\]
 These are the two toric parts of the period lattice.
The complementary Riemann form is principal on the abelian
part. The bounded split approximation also gives the prescribed
\(n'\) which clears the projection to the isotropic line.

Write \(V_{\rm ab}/\Lambda_{\rm ab}\) for the abelian torus in the period decomposition and \(\widetilde y_{\rm ab}\) for the abelian component of a lift of \(y\in Y\). Clearing the projection gives
\[
       n'\widetilde y_{\rm ab}\in\Lambda_{\rm ab}
       \quad\Longrightarrow\quad
       n'[\widetilde y_{\rm ab}]=0\quad\text{in }V_{\rm ab}/\Lambda_{\rm ab}.
\]
Thus the gluing points have the prescribed torsion bound. The rank-one period
comparison keeps the complex flow on this component. The
canonical trivialization and algebraic descent are described
below. The identification of tangent spaces respects the embedding blocks,
so a block change \(X\) of these letters reads the prescribed
split frame.

Changes along the cusp flow multiply that mode by
a nonzero exponential. Thus \eqref{eq:slice-arc-coeff} is
nonzero at one such \(X\), and some monomial of the polynomial
is detected. The field of definition of the split CM data may depend on the coefficient being tested. Taking all such tests gives the rational kernel. After scalar extension, the tests annihilate every complex family whose corresponding low modes vanish.
\label{cor:slice-cm-coefficients-3}

\end{proof}

\paragraph{The normalized degeneration and its fields of definition.}
On the complementary period lattice let \(H_B\) be the Hermitian
form, normalized so that its imaginary part is the integral Riemann
form. Let \(C_y\) be the abelian component of the period belonging
to \(y\in Y\). The toric periods are affine in the arc coordinate
\(w\), with linear term given by the positive form of the preceding
construction. The constant term in their pairing can be chosen to be
\begin{equation}\label{eq:slice-cm-normal-pairing}
 T(y,y')=-\frac{H_B(C_y,C_{y'})}{2i}
                 +\operatorname{Im}H_B(\alpha C_y,C_{y'}),
 \qquad \alpha+\bar\alpha=1,
\end{equation}
where \(\alpha\) is the marked integral element of the CM host.
The Rosati relation and the displayed equation for \(\alpha\)
make \(T\) symmetric. They also give

\[
  T(y,\alpha y')+T(y',\alpha y)=T(y,y').
\]

Consequently \(\exp(2\pi i T(y,\alpha y))\) is the even
trivialization with polarization given by \(T\).
Replacing the old symmetric constant pairing by \(T\) is a
translation along the complex cusp flow. Indeed the difference
is conjugate balanced over \(\mathcal O_E\) on \(Y\), so over
the trace pairing it is a complex linear combination of real
parameters of that flow. The slope and abelian projection are
preserved. At large imaginary \(w\) the positive slope absorbs
the fixed translation in each embedding. A Fourier--Jacobi
coefficient under that translation receives its trace character,
which is never zero.

In the normalized model \(n'C_y\) belongs
to the abelian period lattice. Thus
\begin{equation}\label{eq:slice-cm-tau-normal}
  \exp\bigl(2\pi i T(n'y,n'y')\bigr)
       \exp\bigl(\pi H_B(n'C_y,n'C_{y'})\bigr)=1.
\end{equation}
Indeed its remaining exponent divided by \(2\pi i\) is the
integral Riemann pairing of \(\alpha n'C_y\) and \(n'C_{y'}\).
On this sublattice the two gluing points are the identity of the
abelian variety. The exponential product is the comparison of
the polarization of the trivialization with the rigidification
there. This is the canonical torsion normalization.
\label{cor:slice-cm-arc-model-1}

\begin{proposition}[Comparison of the two evaluations on the normalized CM arc]
\label{prop:marked-cm-arc-base-change}
On the marked CM degeneration with pairing
\eqref{eq:slice-cm-normal-pairing}, let \(s,a_\nu\) and \(s',a'_\nu\) denote the positive slope and the Fourier--Jacobi coefficient before and after the change of constant period. Then
\[
 s'=s>0,\qquad a'_\nu=\gamma_\nu a_\nu,\quad\gamma_\nu\ne0.
\]
 Over the fields with the prescribed marking, the
complex and finite-place evaluations use the same toric parameter and
the same normalized theta and frame lines, up to the indicated constant
frame changes.
\end{proposition}

\begin{proof}
\proofstep{Realizing and descending the marked period.}
The marked CM period realization applies to this complementary
lattice with its principal alternating form, split CM type and
bounded basis. It gives a polarized abelian scheme \(B\) after a
finite number field extension and an identification of its complex
periods and CM action with the above ones. In this construction one
keeps the isometry of marked period lattices, so the same
identification carries the Riemann matrix and the Weil pairings on
the gluing points.

Toroidal uniformization for the isotropic
quotient uses \(B\), the character lattice \(Y\), the gluing
homomorphism \(c:Y\to B[n']\), the positive period form and a
trivialization in the pulled-back cubical line. The norm and the
trivialization are compared using the rigidification at zero and
the symmetric primitive theta characteristic. Writing \(E_B\) for the integral period pairing, the symmetric characteristic evaluates on a period \(a\) as
\[
  (-1)^{E_B(\alpha a,a)}.
\]
 The equality of
the period pairings gives equality of these signs, which is needed
for the two Appell--Humbert readings to agree.
\cite{faltings-chai}\label{cor:slice-cm-arc-model-2}

\proofstep{Preserving the toric parameter in descent.}
To descend the normalized model, first adjoin the \(n'\)-torsion
of \(B\). The characteristic-zero torsion scheme is finite, so
all points of \(c\) are now defined over a number field \(k\).
Choose algebraic generators for the one-dimensional cubical
fibres at \(c(y)\). Compare them with the normalized analytic
generators by nonzero scalars \(q(y)\). In the quotient \(\mathbb C^\times/k^\times\), divide the
analytic and algebraic symmetry and cube identities. They
give an even function \(\bar q\) with bimultiplicative
polarization
\[
 b(y,y')=\frac{\bar q(y+y')}{\bar q(y)\bar q(y')}.
\]
The normalized product in \eqref{eq:slice-cm-tau-normal} is
the rigidification over \(k\) on \(n'\)-multiples, so
\[
 b(y,y')^{(n')^2}=b(n'y,n'y')=1.
\]
Bimultiplicativity also gives \(\bar q(0)=1\).
Since \(b(y,-y)=\bar q(y)^{-2}\), one has

\[
  \bar q(y)^{2(n')^2}=1.
\]
 Thus \(q(y)\) satisfies a
power equation over \(k\) and is algebraic.
For a finite integral generating family of \(Y\), it suffices to
adjoin the scalars on the generators and on their pairwise sums.
In the quotient by the enlarged field all polarizations vanish
by bimultiplicativity, and then all the scalars vanish in the
quotient group by the character law. This is one finite extension
for the entire trivialization. The remaining level data are points
in finite torsion schemes of the bounded level and can be adjoined
in a further finite extension.
If \(k_i/k\) are the finitely many fields of definition of the residual level records, their compositum is finite since
\[
  [k_1\cdots k_t:k]\le\prod_{i=1}^t[k_i:k].
\]
 The rigidified line, the marked
action and the normalized toric parameter are transported together.
\label{cor:slice-cm-arc-model-3}

\paragraph{Finite-place realization of the same arc.}

The archimedean construction has now specified the arc and its field. Its evaluation must also extend to the completed model at every finite place needed for the height.

Use a finite place \(w\) of the field \(K'\) just constructed,
write \(V\) for its complete DVR and \(F_V\) for its fraction
field. Base change the abelian part \(B\) to \(V\). By properness
every gluing point in \(B(F_V)\) extends uniquely to \(B(V)\).
The extended points form a homomorphism \(c_V:Y\to B(V)\):
both sides of 
\[
  c_V(y+y')=c_V(y)+c_V(y')
\]
 extend the same
generic point, as do the sides of the CM equivariance identity.
For the rigidified cubical line \(\mathcal T\) on \(B/V\),
the datum retains each generic generator in
\[
             (c_V(y)^*\mathcal T)\otimes_V F_V .
\]

In particular one retains its valuation relative to
\(c_V(y)^*\mathcal T\), even if the generator is not a basis
over \(V\). This is
the datum used in the integral toroidal model. In particular it
has the same generic fibre as the original selected arc. The
level points in the degeneration datum may be kept as generic
level records. Extension of a generic level point as a point of
an abelian torsion group alone would not specify its toric record.

\proofstep{Integral frames and tangent coordinates on the chart.}
An integral framed chart has constant tangent letters in each
CM embedding block and a generator of the theta line on the
integral abelian part. The comparison of two such descriptions
is made over \(F_V\) by base change of the same degeneration
datum.
Thus two ordered tangent bases \(e,e'\) on a common CM block over the generic fibre differ by
\[
  e'=eA,\qquad A\in{\rm GL}_n(F_V),
\]
 Each block of algebraic tangent letters is related to the
integral block by a constant matrix. Constancy follows by taking
invariant tangent vectors of the semiabelian extension and
identifying their pullbacks via the common generic group. The
theta generators there span the same one-dimensional space.

Let
\(\gamma\) be the scalar expressing the algebraic section in the
integral generator.
Thus if \(\vartheta_{\rm alg}\) and \(\vartheta_{\rm int}\) are the indicated generators, their comparison and its valuation are
\[
  \vartheta_{\rm alg}=\gamma\,\vartheta_{\rm int},\qquad v_w(\vartheta_{\rm alg}/\vartheta_{\rm int})=v_w(\gamma).
\]
 These identifications determine the frame
comparison over the entire arc. A change of the torsion
normal form multiplies the reading by a root of unity, as the
torsion sections differ by their level character. This leaves the
finite order of \(\gamma\) unchanged.
\cite{faltings-chai}\label{cor:slice-cm-arc-model-4}

The generator of the global theta ideal at \(w\) is a local
coordinate for the theta line of this same polarized model.
Carry the generator and a Lie frame through the polarized
pullback. The integral theta generator differs from the former
by a unit in \(V\), since both generate the same rank-one module
and the isomorphism is rigidified. If \(t_w\) generates the theta ideal in this trivialization,
then 
\[
  \gamma^{-1}=u t_w
\]
 with \(u\in V^\times\).

Passing to the dual
ideal changes the sign:
\[
 v_w(\gamma)=-v_w(t_w)=v_w I(\mathcal T^\vee).
\] The determinant
of the Lie comparison is that of the same model. This identifies
the central factor of the coefficient line used at infinity as
well as at finite places. In the finite comparison below, the
integral toroidal evaluation is made on smaller punctured discs
over valued extensions. The input is the framed chart extending
the normalized degeneration, with the valuation-compatible
extension of the coefficient field. Base change makes the two
evaluations agree after the constant frame change. The bound on
the integral reading uses the integral lattice of sections on
that chart.
\label{cor:slice-cm-arc-model-5}
\end{proof}


\paragraph{Finite order of an arc coefficient.}
At a place \(w\) of \(K'\) over \(\lambda\) of \(K\), take
the completed integral semiabelian model. Let
\(\mathcal T\) be the theta line with its normalized generator
and \(\ell_{\tau k}\) the line of the letter in column \(k\)
of the \(\tau\)-block. A toric letter has the unit line.
Write \(\gamma\) for an inverse local generator of the ideal of \(\mathcal T\)
and \(d_{\tau k}\) for the generators of the ideals of the
letter lines. Use \(\operatorname{ord}_w\) with integer orders.
The integral comparison supplies a matrix \(B_\tau\) from
an integral frame to the coordinates normalized by these
generators. In its inverse \(A_\tau\), every entry has order
at least \(-b_w^0\), where
\[
 b_w^0=e_{w/\lambda}
             \max\{\operatorname{ord}_\lambda|\operatorname{disc}E|,
                         \operatorname{ord}_\lambda n'\}.
\]

For a toric letter the infinitesimal character on \(y\in Y\)
is the labelled conjugate of \(p(y)\). A torus tangent vector
over the dual numbers is determined by its evaluations
\(1+\varepsilon h(y)\) on the characters, with \(h\) additive.
The trace-dual condition gives \(d_Ep(y)\in\mathcal O_E\),
so the multiplied tangent vector is integral by this
character description on the same constant torus. For an
abelian letter the local generator \(d_{\tau k}\) gives an
integral tangent vector on \(B\). The classifying extension
has \(n'\)-torsion gluing homomorphism and the chosen
trivialization on \(n'Y\). The pullback under \([n']\) is split,
with these trivializations specifying a splitting compatible
with the generic fibre. Applying its tangent map to the
integral vector on \(B\) extends precisely the generic
semiabelian letter multiplied by \(n'\).

On invariant tangent vectors the abelian multiplication map acts by
\[
  d[n']_0(v)=n'v .
\]
 The local unit
needed to compare a differential model with this frame
is invertible over \(V\). These are separate columns, and
the same bound on the entries uses the larger of the two
orders. The theta generator is the integral rigidified
generator up to a local unit after rescaling by its line.
These comparisons are made on the integral chart of the
toroidal model. An integral section is bounded by one when
evaluated there in any integral frame with unit determinant.
\label{cor:slice-cm-coefficients-4}

\begin{lemma}[Finite and complex bounds for one coefficient]
\label{lem:cm-coefficient-metrics}
Use the theta and tangent lines of the selected arc. For an integral
source form the coefficient bounds in the following calculation hold
on the ideal defined by those lines for the same \(\phi_{\mu,J}\)
at the finite and the complex places. For the finite reading on an integral source form \(F\), the coefficient ideal of the value line gives
\[
 \operatorname{ord}_w\phi_{\mu,J}(F)-\operatorname{ord}_w I(\mathcal L_\mu)\ge-\mathcal A b_w^0.
\]
The complex bound uses \(\max(1,c_1)\) and the frame radius bound \(\Lambda\), independently of the length of the isolating direction.
\end{lemma}

\begin{proof}
\proofstep{The nonarchimedean coefficient estimate.}
We give the coefficient argument because it also fixes the finite
offset in the quotient. In the variable
\(Y=(Y_\tau)\) for the integral frames put
\[
 R(Y,t)=\gamma^{-n}P_F((Y_\tau B_\tau D_\tau^{-1})_\tau,t),
 \qquad D_\tau=\operatorname{diag}(d_{\tau k}^{-1}),
 \qquad A_\tau=B_\tau^{-1}.
\]
The finite chart comparison gives convergence and
\(|R(Y,t)|_w\le1\) for the integral \(Y\) with unit block
determinants and the allowed \(t\) with \(0<|t|_w<1\).
One may evaluate over the completed algebraic closure with
infinite residue field. The allowed radii are dense in
\((0,1)\) and the readings of \(t\) are stable under multiplication
by units with arbitrary nonzero residues in an infinite
subfield.
In particular unit variation on any allowed circle preserves its radius:
\[
  |t_0u|_w=|t_0|_w\qquad(u\in V^\times).
\]
 Such extensions and parameters in the integral arc
comparison leave the value line on \(K'\) fixed.

For a fixed \(Y\) and allowed radius \(|t_0|_w\), the maximum
of \(\big|[t^i]R(Y,t)t_0^i\big|_w\) is attained if it is nonzero.
Convergence makes only finitely many of the terms attain it.
Divide by one such term and reduce the sum at \(t=t_0u\).
If \(a_i=[t^i]R(Y,t)\), the normalization is by one term attaining
\[
  |a_jt_0^j|_w=\max_i |a_it_0^i|_w>0 .
\]
The remaining terms of maximal norm give a nonzero Laurent
polynomial in the residue of \(u\). A unit with a residue
outside its finitely many roots gives a value of precisely
the maximal norm. Hence
\[
 \big|[t^J]R(Y,t)\big|_w\,|t_0|_w^J\le1.
\]
Let the radii tend to one from below. This works for every
integer \(J\).

\proofstep{Clearing denominators of the chosen coefficient at finite places.}
For the fixed polynomial \([t^J]R(Y,t)\) divide again by a
coefficient of maximal absolute value. Its reduction times
the product of the block determinants is nonzero. A tuple
in the infinite residue subfield detects this polynomial,
so lifts of the tuple give unit block determinants and
detect the original maximal coefficient norm. It follows
that every coefficient of \([t^J]R(Y,t)\) is integral.
The reverse substitution 
\[
  Y_\tau=X_\tau D_\tau A_\tau
\]
uses \(\mathcal A\) factors from the matrices \(A_\tau\)
in every term. Their contribution to a coefficient has
order at least \(-\mathcal A b_w^0\), with no additional
loss from the integer multinomial coefficients or the sums.
Evaluation after the remaining diagonal rescaling multiplies
the \(\mu\)-coefficient by exactly
\(\prod_{\tau,i,k}d_{\tau k}^{-\mu_{\tau i k}}\).

Consequently the finite estimate is
\[
 \operatorname{ord}_w\phi_{\mu,J}(F)-
       \operatorname{ord}_w I(\mathcal L_\mu)
       \ge-\mathcal A b_w^0,
 \qquad
 \mathcal L_\mu=(\mathcal T^\vee)^n
                  \otimes\bigotimes_{\tau,i,k}
                         (\ell_{\tau k}^\vee)^{\mu_{\tau i k}}.
\]
This gives \(s_{iw}=-\operatorname{ord}_w I(\mathcal L_\mu)\)
in the quotient calculation below. A test of the wrong
total frame degree vanishes by homogeneity. For the other
tests the common baseline has the factor \(\mathcal A\).
\label{cor:slice-cm-coefficients-5}

\paragraph{Archimedean bound for the same coefficient.}

The local order accounts for the finite part of the height. The archimedean contribution must be evaluated on the same coefficient and the same arcs.

Use a good embedding \(\sigma\) of \(K'\) and denote the
positive metric factors of \(\mathcal T,\ell_{\tau k}\) there
by \(q_\theta,q_{\tau k}\). Put \(v=\max(1,c_1)\) and
\(\epsilon_*=v^{-1}\). Move the frame in
\eqref{eq:slice-arc-coeff} along its cusp flow to all CM
times equal to \(\epsilon_*\). The shift adds
\(\epsilon_*-y\tau(s_E)\) to each old time. The Fourier
factor at time \(y\tau(s_E)\) is recovered by multiplying by
\[
 \exp\left(4\pi\sum_{\tau}\tau(\nu_0)
                         (\epsilon_*-y\tau(s_E))\right)
 = |t|^J\exp\left(4\pi\epsilon_*\sum_\tau\tau(\nu_0)\right).
\]
Here \(y=\operatorname{Im}w\) and the sum is over the CM half.

The total frequency bound and the common shifted time give
\[
  \epsilon_*\sum_\tau\tau(\nu_0)\le\frac{c_1}{\max(1,c_1)}\le1,\qquad \exp\!\left(4\pi\epsilon_*\sum_\tau\tau(\nu_0)\right)\le e^{4\pi}.
\]
 The \(|t|^J\) cancels the factor
from taking the modulus of \eqref{eq:slice-arc-coeff}.
Thus the initial height and the size of the isolating direction
\(s_E\) add no cost to this estimate. In particular the proof of
isolation needs no uniform bound on \(R\). The transvection torus integral for the coefficient
at the new time bounds it by the supremum of the section
divided by its pointwise normalizing factors.

\proofstep{Archimedean normalization on the same arc.}
For clarity, the remaining normalizing factors in this
integral can be computed before coefficient extraction in
\(X\). Let \(H_X\) be the positive Riemann form in the
shifted frame and \(h_{\tau i}=\pi (H_X)_{\tau i,\tau i}\).
The complex comparison of the theta and frame
normalizations gives the factor
\[
 \left(C v^f\right)^{n/2}\prod_{\tau,i}h_{\tau i}^{a_{\tau i}/2},
 \qquad C\le q_\theta^2\operatorname{rd}(E)^f .
\]
Indeed the central normalization is initially
\((C\prod_\tau \mathrm{time}_\tau^{-1})^{n/2}\)
times the central-weight norm in the old frame.

The central-weight norm times this product is invariant
under the move. Dividing by the norm of the total weight
at the new frame leaves the diagonal factors displayed
above. In the normalized theta trivialization, comparison
with the period volume gives \(C/q_\theta^2\le
\sqrt{|\operatorname{disc}E|}\).
These comparisons use the Riemann form on the same
semiabelian degeneration as the trace formula and the
abelian CM metric on its central fibre.
\label{cor:slice-cm-coefficients-6}

\proofstep{Bounding the coefficient on the polydisc.}
Give \(X_{\tau i k}\) radius \(q_{\tau k}^{-1}\) in a closed
complex polydisc. At new time \(\epsilon_*\), the individual
CM letter satisfies

\[
  \sqrt{\pi H_X(e_{\tau k},e_{\tau k})}\le
\Lambda\sqrt v\,q_{\tau k},
\]
 where the pairing on the left
is the pullback of the shifted Riemann form to the letter.
The same inequality also holds for the toric letter with
its unit line; its squared length is bounded by the
reciprocal time. The row is a sum of at most \(r\) such
letters with the indicated coefficients. The triangle
inequality in that Hermitian norm gives
\(\sqrt{h_{\tau i}}\le r\Lambda\sqrt v\).
Put
\[
 p_{\rm arch}=\log(r\Lambda)
       +\tfrac12\log v+\log\operatorname{rd}(E).
\]
The product above is at most
\(q_\theta^n\exp(p_{\rm arch}(\mathcal A+fn))\).
These bounds first apply where the block matrices are
invertible.

Since the reading is a polynomial in \(X\),
continuity extends the bound to the full polydisc.
The multivariate Cauchy coefficient formula gives
\[
 |\sigma\phi_{\mu,J}(F)|\,\|\mathcal L_\mu\|_\sigma
 \le \|F\|_{\sup,\sigma|K}\,
       \exp\big(4\pi+p_{\rm arch}(\mathcal A+fn)\big).
\]
In fact the product of the Cauchy radii is exactly the
differential metric factor of the dual value line. The
division by \(q_\theta^n\) gives its remaining factor.
At the other embeddings use the conjugate good
uniformization, with the same metric by complex
conjugation. This gives the archimedean offset for precisely
the ideal in the finite calculation.
\label{cor:slice-cm-coefficients-7}
\end{proof}
